\section{Background}
\label{s:back}

\subsection{The Internal $\pi$-calculus}
\label{ss:piI}

The Internal $\pi$-calculus, $\piI$, is, intuitively, 
a subset of the $\pi$-calculus 
in which all outputs are bound. This is syntactically enforced by
having outputs written as $\bout a\tilb$ (which in the $\pi$-calculus would be an
abbreviation for $\res\tilb \out a \tilb$). 
We use the polyadic version of the calculus, in which a tuple of names may be exchanged in
a communication step. 
All tuples of names in \piI 
are made of pairwise distinct components. 

\emph{Abstractions} are used to write name-parametrised processes,  for
instance, when writing recursive process definitions.
 The instantiation of the parameters of an
abstraction $B$ is done via the {\em application} construct $\app B \tila$.  
Processes and abstractions form the set of \emph{agents}, ranged over by $\AG$.
Lowercase letters
$a,b, \ldots, x,y, \ldots$ range over the infinite set of names.  The grammar of \piI
 is thus:
\[\begin{array}{ccll}
  P & \defeq & \nil    \midd    \inp a \tilb . P    \midd    \bout a \tilb . P 
    \midd    \res a P \midd   P_1 |  P_2   \midd  ! \inp a \tilb . P   
  \midd \app  B \tila & \mbox{(processes)} \\[\myptD]
  B & \defeq & \bind \tila P \midd {\rmmm K} & \mbox{(abstractions)}
\end{array}\]
We omit the standard definition of free names,
bound names, and names of an agent, respectively indicated with $\fn -$, $\bn -$, and
$\n -$. The grammar for processes includes the operators  
$\nil$ (representing  the inactive process), 
input and output prefixes, restriction, parallel composition, replication, and
application.   

In the grammar of abstractions, 
${\rmmm K}$ is a constant,  used to write recursive definitions. Each
constant ${\rmmm K}$ has a defining equation of the form 
${\rmmm K} \defeq  \bind{\tilx} P$, where $\bind{\tilx} P  $ is name-closed (that is, 
without free names); $\tilx$ 
 are the formal parameters of
the constant (replaced by the actual parameters whenever the constant
is used).
Replication could be avoided in the syntax since it can be encoded
with recursion. However its semantics is simple, and it is 
a useful construct for encodings, 
thus we chose to include it in the grammar.

\ifcutLMCS
Sometimes, we use $\defi$ as an abbreviation mechanism, to
assign a name to an expression to which we want to refer later.
\gj{Is this use of $\defi$ used anywhere?}
\fi


% \paragraph{Abstraction and application.}
An {application redex}, that is, a term syntactically of the form 
 $\app{((\til x) P)}{\til a}$, can be normalised as $P \subst {\til x}{\til a}$. An agent is
\emph{normalised} if all such application redexes have been
contracted.
Although the full
application $\app  {\rmmm K}{\til a} $ is often convenient, when
it comes to reasoning on behaviours 
it is useful to assume that all expressions are normalised, in the above sense.
Thus in the remainder of the paper 
\emph{we identify an agent with its normalised expression}.
\reviewComment{does this include recursive definitions? In
which case, are normalized expressions infinite?}
\reviewCommentDS{Answer: no, only expressions of that syntactic form are
  modified. Recursive definition are only unfolded via rule {\trans{con}} pf Figure 1. This
kind of normalisation is a (finite) syntactic simplification; it is not mandatory, however
it sometimes simplifies  structural reasoning  on processes. We have rephrased a bit the sentence.}

Since the calculus is polyadic, 
we assume a \emph{sorting system}~\cite{Mil91}
   to avoid disagreements  in the arities of
the tuples of names 
carried by a given name 
and in applications of abstractions.
Being not essential, it will not be presented here.
The reader should
take for granted that all agents described obey a sorting. 

% Substitutions are of the form $\sub \tilb \tila$, and
% are   finite assignments of  names to names.
% We use $\sigma$ and $\rho$  to range over substitutions. 
% The application of a
%   substitution  $\sigma $ to an expression $H$ is
% written $H \sigma $.
% Substitutions have 
% precedence over the operators of the language; $\sigma  
% \rho$ is the composition of substitutions  where $\sigma $ is performed
% first, 
%  therefore $P \sigma  
% \rho$ is $ (P \sigma ) \rho$.


\subsubsection{Operational semantics and behavioural relations}
\label{sss:os_piI}

In Figure~\ref{t:pii_ST}
we report the standard LTS of \piI.
The symbol $\equival$ is used to indicate $\alpha$-equivalence
between \piI agents.
Rule  {\trans{abs}} is used to formalise, by means of an action, the ground instantiation
of the parameters of an abstraction. 
In both tables, we omit the symmetric of {\tt parL}, called {\tt parR}.
In rules {\tt com}, $\outC \mu$ is the complementary of $\mu$, 
thus if $\mu = \inp  a \tilx$
then  $\outC\mu = \bout  a \tilx$, and conversely if  $\mu = \bout  a \tilx$
then  $\outC\mu = \inp  a \tilx$. 


\begin{figure*}[ht]
\begin{center}
\begin{tabular}{@{}r@{\,}l@{\hspace{1.2em}}r@{\,}l@{\hspace{1.2em}}r@{\,}l@{}}
{\trans{alpha}}: &  $\displaystyle{  P\equival P' \andalso   { P'} \arrPI\mu    { P''}
  \over 
  P \arrPI\mu {P''}  
 } $ &  
{\trans{pre}}: &  
$\displaystyle{%\bn \mu = \tilx \andalso \tilx \cap \phi = \emptyset  
\over
{\mu.      P} \arrPI\mu {P}}$   
& {\trans{parL}}: &  $\displaystyle{   P \arrPI\mu { P'} \over  { P | Q}   \arrPI\mu
{P'| Q} } $   
\\[\mysp] 
{\trans{res}}: & $\displaystyle{{  P} \arrPI{\mu}{P'}
\andalso  x \not \in \n\mu
 \over
{ \res x     P}   \arrPI{ \mu}{\res x P'  }} $ 
&
\multicolumn{4}{c}{ 
    {\trans{com}}:  \; \;    $\displaystyle{ {P} \arrPI{ \mu }{P'} \andalso   
 Q
\arrPI{\outC\mu}{Q'}
\andalso \mu\neq \tau,   \; \tilx = \bn\mu 
  \over
{     P |  Q }\arrPI{ \tau} 
{
 \res{  \tilx} (P'
|  Q') } }$    
}
 \\[\mysp]
    {\trans{rep}}: &
 $\displaystyle{ {\mu.  P }\arrPI{\mu} {P'}
 \over
{ !\mu. P }   \arrPI{ \mu}{ P' | !\mu. P } } $  
&
\multicolumn{4}{c}{ 
    {\trans{con}}: \;\;
$  \displaystyle{   P \arrPI \mu {P'}
\andalso  \mbox{ } {\rmmm K}  \defi
  \abs{\til y} Q  
\andalso 
\mbox{ } \abs{\til y} Q \equival \abs{\til x} P
    \over 
{
 \app {\rmmm K}{\til x}}  \arrPI\mu 
{P'} }$
}
\\[\mysp] 
\multicolumn{6}{c}{ 
    {\trans{abs}}:  \; \;    
$  \displaystyle{ 
\mbox{( $B =  \abs{\til y} Q $   or 
($B = K $ and $K \defi   \abs{\til y} Q $) )}  
  \andalso  \abs{\til y} Q \equival \abs{\til x} P
  \over 
{B \arr{({\til x})}     P }
 }$
}
\end{tabular} 
\end{center} 
\caption{The standard LTS for \piI}
\label{t:pii_ST}
 \Mybar 
\end{figure*} 


In the LTS for \piI,
transitions are of the form $\AG \arrPI \mu \AG '$, where the bound names of $\mu$ are {fresh},
i.e., they  do not appear free
in $\AG$.  The choice of bound names can be made more explicit by following the operational game
semantics approach in which transitions are defined on configurations that include a name
set component. 
(Some presentations of $\pi$-calculi indeed make use of configurations.)  
The two approaches are interchangeable, and the choice is largely a matter of taste. 
The standard LTS  is however more convenient in reasoning, e.g.,  the up-to
proof techniques of Section~\ref{ss:upto}.

All behavioural relations are extended to abstractions by requiring ground
instantiation of the parameters; this is expressed by means of a transition; e.g., the action 
$ \bind{\tilx} P \arrPI {\bind{\tilx}{}}  P$ (see rule  {\trans{abs}} in Figure~\ref{t:pii_ST}). 

\finish{see if we want to mention also structural congruence $\equiv$} 

\subsection{The expansion relation}


The expansion relation, $\expa$~\cite{SaWabook}, 
is an asymmetric variant of 
$\wbPI$  which allows us to  count the number of $\tau$-actions
performed by the processes. 
Thus, $P \expa Q$ holds if $P \wbPI
Q$ but also $Q $ has at least as many $\tau$-moves as $P$. 
We write $P \arcap\mu P'$  
when either $P \arr\mu P'$ or ($\mu =\tau$ and
$P=P'$).

\begin{definition}[expansion]
  A relation $\R \subseteq   \pr \times \pr $ 
  is an {\em expansion } if $P \RR Q$
  implies:
  \begin{enumerate}
  \item  Whenever 
  $P \arr\mu P'$,    
  there exists  $Q' $  s.t.\ $Q \ArrPI\mu Q'$ 
  and $P'  \RR  Q'$;
  %%
  \item whenever $Q \arr\mu Q'$, 
  there exists  $P' $  s.t.\ $P \arcap\mu P'$ 
  and $P'  \RR  Q'$.
  \end{enumerate}
  \noindent
  We say that $Q$  {\em expands } $P$, written
  $P \expa  Q$, 
  if $P  \R  Q$,  for some expansion $ \R $.   
\end{definition} 


Relation $\expa$ is a precongruence, and  is
strictly included in $\wbPI$.  
